\makeatletter
\expandafter\gdef\csname CGTip@sec:introduction\endcsname{Section 1: Introduction (p. 2)}
\expandafter\gdef\csname CGTip@subsec:object\endcsname{Section 1.1: The object (p. 2)}
\expandafter\gdef\csname CGTip@eq:intro_clock\endcsname{Eq. (1.1), p. 2: tau\_k=(2pi (k-(11)/(8)))/(W((k-(11)/(8))/e))(k>=1),}
\expandafter\gdef\csname CGTip@eq:intro_FL_equation\endcsname{Eq. (1.2), p. 2: (t)/(2pi)log(t)/(2pi e)+(1)/(pi)lim\_eps ->0\textasciicircum{}+arg zeta ((1)/(2)+eps+it)=k-(11)/(8)}
\expandafter\gdef\csname CGTip@eq:intro_defect\endcsname{Eq. (1.3), p. 2: u\_k(X)=k-(1)/(2)-N\_0(gamma\_k)-S\_X(gamma\_k),}
\expandafter\gdef\csname CGTip@subsec:why\endcsname{Section 1.2: Why it is interesting (p. 3)}
\expandafter\gdef\csname CGTip@subsec:main_results\endcsname{Section 1.3: Main results (p. 3)}
\expandafter\gdef\csname CGTip@eq:intro_main\endcsname{Eq. (1.4), p. 3: (1)/(no.I\_T)sum\_kin I\_Tu\_k(X)\textasciicircum{}2=(1)/(2pi\textasciicircum{}2)log(1)/(c)+O(1+sqrt(log(1/c))).}
\expandafter\gdef\csname CGTip@subsec:selberg_relation\endcsname{Section 1.4: Relation to Selberg's work (p. 3)}
\expandafter\gdef\csname CGTip@sec:notation\endcsname{Section 2: Notation and the clock (p. 3)}
\expandafter\gdef\csname CGTip@subsec:counting\endcsname{Section 2.1: Counting functions (p. 3)}
\expandafter\gdef\csname CGTip@eq:RvM\endcsname{Eq. (2.1), p. 3: N(t)=(1)/(pi)vartheta(t)+1+S(t), S(t)=(1)/(pi)arg zeta ((1)/(2)+it)=O(log t),}
\expandafter\gdef\csname CGTip@eq:N0_def\endcsname{Eq. (2.2), p. 3: Theta\_0(t)=(t)/(2)log(t)/(2pi e)-(pi)/(8), N\_0(t)=(1)/(pi)Theta\_0(t)+1, omega(t)=vartheta(t)-Theta\_0(t),}
\expandafter\gdef\csname CGTip@eq:N_decomposition\endcsname{Eq. (2.3), p. 3: N(t)=N\_0(t)+S(t)+(omega(t))/(pi) .}
\expandafter\gdef\csname CGTip@eq:L_def\endcsname{Eq. (2.4), p. 4: L(t)=N\_0'(t)=(1)/(2pi)log(t)/(2pi).}
\expandafter\gdef\csname CGTip@eq:IT_size\endcsname{Eq. (2.5), p. 4: no.I\_T=TL(T)+O(T).}
\expandafter\gdef\csname CGTip@subsec:clock\endcsname{Section 2.2: The clock and the displacements (p. 4)}
\expandafter\gdef\csname CGTip@lem:clock\endcsname{Lemma 2.1 (p. 4): The sequence (1.1) is the unique solution of N\_0(tau\_k)=k-(1)/(2) with tau\_k>2pi, k>=2. As a function of a real variable k, tau\_k is smooth and increasing with (dtau\_k)/(dk)=(1)/(L(tau\_k))=(2pi)/(log(tau\_k/2pi)).}
\expandafter\gdef\csname CGTip@eq:clock_derivative\endcsname{Eq. (2.6), p. 4: (dtau\_k)/(dk)=(1)/(L(tau\_k))=(2pi)/(log(tau\_k/2pi)).}
\expandafter\gdef\csname CGTip@def:displacements\endcsname{Definition 2.2 (p. 4): The ordinal displacement and the displacement of the clock are s\_k=k-(1)/(2)-N\_0(gamma\_k), delta\_k=tau\_k-gamma\_k .}
\expandafter\gdef\csname CGTip@eq:displacements\endcsname{Eq. (2.7), p. 4: s\_k=k-(1)/(2)-N\_0(gamma\_k), delta\_k=tau\_k-gamma\_k .}
\expandafter\gdef\csname CGTip@eq:s_delta\endcsname{Eq. (2.8), p. 4: s\_k=L(xi\_k) delta\_kfor some xi\_k between gamma\_k and tau\_k .}
\expandafter\gdef\csname CGTip@rem:gram_points\endcsname{Remark 2.3 (p. 4): Since Theta\_0(tau\_k)=pi(k-(3)/(2)) and vartheta=Theta\_0+omega, the clock point tau\_k is the point at which vartheta takes the value pi(k-(3)/(2)), up to a shift O(tau\_k\textasciicircum{}-1(log tau\_k)\textasciicircum{}-1). Thus tau\_k lies, in phase, half way between the Gram points g\_k-2 and g\_k-1, defined by vartheta(g\_n)=npi. The integer displaceme...}
\expandafter\gdef\csname CGTip@lem:s_bounds\endcsname{Lemma 2.4 (p. 4): For every k with gamma\_k>2, N(gamma\_k\textasciicircum{}-)-N\_0(gamma\_k)+(1)/(2)<= s\_k<= N(gamma\_k\textasciicircum{}+)-N\_0(gamma\_k)-(1)/(2), where N(gamma\textasciicircum{}+-) are the one-sided limits. In particular |s\_k|<=|S(gamma\_k\textasciicircum{}-)|+|S(gamma\_k\textasciicircum{}+)|+1.}
\expandafter\gdef\csname CGTip@eq:s_bounds\endcsname{Eq. (2.9), p. 4: N(gamma\_k\textasciicircum{}-)-N\_0(gamma\_k)+(1)/(2)<= s\_k<= N(gamma\_k\textasciicircum{}+)-N\_0(gamma\_k)-(1)/(2),}
\expandafter\gdef\csname CGTip@subsec:euler_phase\endcsname{Section 2.3: Euler phases and the truncated equation (p. 4)}
\expandafter\gdef\csname CGTip@eq:an_def\endcsname{Eq. (2.10), p. 4: a\_n=-(Lambda(n))/(pisqrt(n) log n),}
\expandafter\gdef\csname CGTip@eq:SA_def\endcsname{Eq. (2.11), p. 4: S\_A(t)=sum\_nin Aa\_nsin(tlog n), v\_A=(1)/(2)sum\_nin Aa\_n\textasciicircum{}2 .}
\expandafter\gdef\csname CGTip@eq:SX_def\endcsname{Eq. (2.12), p. 4: S\_X(t)=-(1)/(pi)sum\_p<= Xarg (1-p\textasciicircum{}-1/2-it)=-(1)/(pi)sum\_p<= Xsum\_m>=1(sin(mtlog p))/(m p\textasciicircum{}m/2),}
\expandafter\gdef\csname CGTip@eq:cP_def\endcsname{Eq. (2.13), p. 5: c\_P=(1)/(2pi\textasciicircum{}2)sum\_pin PLi\_2((1)/(p))=lim\_M->infv\_p\textasciicircum{}m: pin P, m<= M,}
\expandafter\gdef\csname CGTip@eq:sigma_M_def\endcsname{Eq. (2.14), p. 5: |S\_X(t)|<=sigma\_X:=(1)/(pi)sum\_p<= Xsum\_m>=1(1)/(m p\textasciicircum{}m/2), |S\_X'(t)|<= M\_X:=(1)/(pi)sum\_p<= X(log p)/(sqrt(p)-1),}
\expandafter\gdef\csname CGTip@def:LeClair_model\endcsname{Definition 2.5 (p. 5): For X>=2 the truncated counting model and the defect of the k-th zero are F\_X(t)=N\_0(t)+S\_X(t), u\_k(X)=k-(1)/(2)-F\_X(gamma\_k)=s\_k-S\_X(gamma\_k). The one-step corrected clock is tau\_k\textasciicircum{}X=tau\_k-(S\_X(tau\_k))/(L(tau\_k))=tau\_k-(dtau\_k)/(dk) S\_X(tau\_k), and, for a saturation level B>0, tau\_k\textasciicircum{}X,B=tau\_k-sat\_B(S\_X(tau\_k)/L(tau...}
\expandafter\gdef\csname CGTip@eq:defect_def\endcsname{Eq. (2.15), p. 5: F\_X(t)=N\_0(t)+S\_X(t), u\_k(X)=k-(1)/(2)-F\_X(gamma\_k)=s\_k-S\_X(gamma\_k).}
\expandafter\gdef\csname CGTip@eq:one_step_clock\endcsname{Eq. (2.16), p. 5: tau\_k\textasciicircum{}X=tau\_k-(S\_X(tau\_k))/(L(tau\_k))=tau\_k-(dtau\_k)/(dk) S\_X(tau\_k),}
\expandafter\gdef\csname CGTip@lem:LeClair_equation\endcsname{Lemma 2.6 (p. 5): Equation (1.2) is equivalent to N\_0(t)+S(t)=k-(1)/(2). Replacing (1)/(pi)arg zeta((1)/(2)+it) by (1)/(pi)arg prod\_p<= X(1-p\textasciicircum{}-1/2-it)\textasciicircum{}-1 turns it into F\_X(t)=k-(1)/(2), and tau\_k\textasciicircum{}X is the first step of the iteration t-> t-(F\_X(t)-k+(1)/(2))/L(t) started at t=tau\_k.}
\expandafter\gdef\csname CGTip@sec:known_results\endcsname{Section 3: Known results used (p. 5)}
\expandafter\gdef\csname CGTip@thm:known_basic\endcsname{Theorem 3.1 (Riemann--von Mangoldt; Littlewood; [ref, Thms. 9.4, 9.9 and S 9.4], [ref]) (p. 5): S(t)=O(log t); the gaps gamma\_k+1-gamma\_k are bounded; and S\_1(t)=int\_0\textasciicircum{}tS(u) du=O(log t).}
\expandafter\gdef\csname CGTip@thm:known_BW\endcsname{Theorem 3.2 (Bellotti--Wong [ref]) (p. 5): For T>= e, |N(T)-(T)/(2pi)log(T)/(2pi e)|<=0.10076log T+0.24460log log T+8.08344 .}
\expandafter\gdef\csname CGTip@eq:BW_bound\endcsname{Eq. (3.1), p. 5: |N(T)-(T)/(2pi)log(T)/(2pi e)|<=0.10076log T+0.24460log log T+8.08344 .}
\expandafter\gdef\csname CGTip@thm:known_selberg_moments\endcsname{Theorem 3.3 (Selberg [ref]; see [ref, p. 225] and [ref, S 9.25]) (p. 5): For every integer j>=1, with c\_j=((2j)!)/(j! (2pi)\textasciicircum{}2j), int\_T\textasciicircum{}2TS(t)\textasciicircum{}2j dt=c\_j T(log log T)\textasciicircum{}j+O\_j(T(log log T)\textasciicircum{}j-1/2).}
\expandafter\gdef\csname CGTip@eq:selberg_moments\endcsname{Eq. (3.2), p. 5: int\_T\textasciicircum{}2TS(t)\textasciicircum{}2j dt=c\_j T(log log T)\textasciicircum{}j+O\_j(T(log log T)\textasciicircum{}j-1/2).}
\expandafter\gdef\csname CGTip@thm:known_selberg_approx\endcsname{Theorem 3.4 (Selberg's approximation, in the form of [ref, S 3, pp. 232--233]; [ref]) (p. 5): Let x\_0=T\textasciicircum{}1/160 and S\textasciicircum{}rm pr\_x\_0(t)=-(1)/(pi)sum\_p<x\_0p\textasciicircum{}-1/2sin(tlog p). Then int\_T\textasciicircum{}2T(S(t)-S\textasciicircum{}rm pr\_x\_0(t))\textasciicircum{}4dt<< T .}
\expandafter\gdef\csname CGTip@eq:selberg_approx\endcsname{Eq. (3.3), p. 5: int\_T\textasciicircum{}2T(S(t)-S\textasciicircum{}rm pr\_x\_0(t))\textasciicircum{}4dt<< T .}
\expandafter\gdef\csname CGTip@thm:known_MV\endcsname{Theorem 3.5 (Montgomery--Vaughan [ref, Cor. 3]) (p. 6): For complex b\_n, int\_T\textasciicircum{}2T|sum\_n<= Yb\_nn\textasciicircum{}-it|\textasciicircum{}2dt=sum\_n<= Y|b\_n|\textasciicircum{}2(T+O(n)) with an absolute implied constant.}
\expandafter\gdef\csname CGTip@thm:known_mertens\endcsname{Theorem 3.6 (Mertens [ref]) (p. 6): For y>=2, sum\_p<= y(log p)/(p)=log y+O(1) and sum\_p<= y(1)/(p)=log log y+B\_1+O((1)/(log y)) with a constant B\_1.}
\expandafter\gdef\csname CGTip@thm:known_KL\endcsname{Theorem 3.7 (Kusmin--Landau; [ref, Thm. 2.1]) (p. 6): Let f be real with a continuous monotone derivative on [a,b], and suppose that ||f'(x)||>=lambda>0 there, where ||*|| is the distance to the nearest integer. Then sum\_a<n<= be\textasciicircum{}2pi if(n)<<lambda\textasciicircum{}-1.}
\expandafter\gdef\csname CGTip@thm:known_landau\endcsname{Theorem 3.8 (Landau [ref]; uniform version by Gonek [ref]) (p. 6): For fixed x>1, sum\_0<gamma <= Tx\textasciicircum{}rho=-(T)/(2pi)Lambda(x)+O(log T), where Lambda(x)=0 unless x is a prime power, and the zeros are counted with multiplicity.}
\expandafter\gdef\csname CGTip@thm:known_RH_S\endcsname{Theorem 3.9 (Littlewood; [ref, Thm. 14.13]) (p. 6): On the Riemann hypothesis, S(t)=O(log t/log log t).}
\expandafter\gdef\csname CGTip@sec:identities\endcsname{Section 4: Exact identities at the zeros (p. 6)}
\expandafter\gdef\csname CGTip@thm:identity_general\endcsname{Theorem 4.1 (p. 6): Let m<n be indices, let F be a real C\textasciicircum{}1 function on [gamma\_m,gamma\_n] and put u\_k=k-(1)/(2)-F(gamma\_k) for m<= k<= n. For every integer j>=1, int\_gamma\_m\textasciicircum{}gamma\_n(N(t)-F(t))\textasciicircum{}2jF'(t) dt=sum\_k=m+1\textasciicircum{}n-1P\_j(u\_k)+(((1)/(2)+u\_m)\textasciicircum{}2j+1+((1)/(2)-u\_n)\textasciicircum{}2j+1)/(2j+1), where P\_j(u)=(((1)/(2)+u)\textasciicircum{}2j+1+((1)/(2)-u)\textasciicircum{}2j+1)/(2j+1)=sum\_i=0...}
\expandafter\gdef\csname CGTip@eq:identity_general\endcsname{Eq. (4.1), p. 6: int\_gamma\_m\textasciicircum{}gamma\_n(N(t)-F(t))\textasciicircum{}2jF'(t) dt=sum\_k=m+1\textasciicircum{}n-1P\_j(u\_k)+(((1)/(2)+u\_m)\textasciicircum{}2j+1+((1)/(2)-u\_n)\textasciicircum{}2j+1)/(2j+1),}
\expandafter\gdef\csname CGTip@eq:Pj_def\endcsname{Eq. (4.2), p. 6: P\_j(u)=(((1)/(2)+u)\textasciicircum{}2j+1+((1)/(2)-u)\textasciicircum{}2j+1)/(2j+1)=sum\_i=0\textasciicircum{}jbinom(2j+1,2i)(2\textasciicircum{}2i-2j)/(2j+1) u\textasciicircum{}2i=u\textasciicircum{}2j+sum\_i<jb\_j,iu\textasciicircum{}2i.}
\expandafter\gdef\csname CGTip@eq:identity_quadratic\endcsname{Eq. (4.3), p. 6: sum\_k=m+1\textasciicircum{}n-1u\_k\textasciicircum{}2=int\_gamma\_m\textasciicircum{}gamma\_n(N(t)-F(t))\textasciicircum{}2F'(t) dt-(n-m-1)/(12)-(((1)/(2)+u\_m)\textasciicircum{}3+((1)/(2)-u\_n)\textasciicircum{}3)/(3) .}
\expandafter\gdef\csname CGTip@eq:identity_one_gap\endcsname{Eq. (4.4), p. 6: int\_gamma\_k\textasciicircum{}gamma\_k+1(k-F(t))\textasciicircum{}2jF'(t) dt=int\_F(gamma\_k)\textasciicircum{}F(gamma\_k+1)(k-y)\textasciicircum{}2jdy=(((1)/(2)-u\_k+1)\textasciicircum{}2j+1+((1)/(2)+u\_k)\textasciicircum{}2j+1)/(2j+1),}
\expandafter\gdef\csname CGTip@rem:twelfth\endcsname{Remark 4.2 (p. 7): The constant (1)/(12) is the mean square of a centred sawtooth of unit step: between two consecutive zeros N-F decreases through one unit in the variable y=F(t). Identity (4.3) therefore says that the mean square of the defect at the zeros is the dF-mean square of N-F minus the contribution of the sawtooth. Nothing...}
\expandafter\gdef\csname CGTip@cor:identity_clock\endcsname{Corollary 4.3 (p. 7): With F=N\_0, so that u\_k=s\_k and N-N\_0=S+omega/pi almost everywhere, sum\_k=m+1\textasciicircum{}n-1s\_k\textasciicircum{}2=int\_gamma\_m\textasciicircum{}gamma\_n(S(t)+(omega(t))/(pi))\textasciicircum{}2L(t) dt-(n-m-1)/(12)-(((1)/(2)+s\_m)\textasciicircum{}3+((1)/(2)-s\_n)\textasciicircum{}3)/(3) . With F=F\_X, so that u\_k=u\_k(X) and N-F\_X=S-S\_X+omega/pi almost everywhere, sum\_k=m+1\textasciicircum{}n-1u\_k(X)\textasciicircum{}2=int\_gamma\_m\textasciicircum{}gamma\_n(S(t)-S\_X(...}
\expandafter\gdef\csname CGTip@eq:identity_clock\endcsname{Eq. (4.5), p. 7: sum\_k=m+1\textasciicircum{}n-1s\_k\textasciicircum{}2=int\_gamma\_m\textasciicircum{}gamma\_n(S(t)+(omega(t))/(pi))\textasciicircum{}2L(t) dt-(n-m-1)/(12)-(((1)/(2)+s\_m)\textasciicircum{}3+((1)/(2)-s\_n)\textasciicircum{}3)/(3) .}
\expandafter\gdef\csname CGTip@eq:identity_LeClair\endcsname{Eq. (4.6), p. 7: sum\_k=m+1\textasciicircum{}n-1u\_k(X)\textasciicircum{}2=int\_gamma\_m\textasciicircum{}gamma\_n(S(t)-S\_X(t)+(omega(t))/(pi))\textasciicircum{}2(L(t)+S\_X'(t))dt-(n-m-1)/(12)-(((1)/(2)+u\_m)\textasciicircum{}3+((1)/(2)-u\_n)\textasciicircum{}3)/(3) .}
\expandafter\gdef\csname CGTip@lem:first_moment_identity\endcsname{Lemma 4.4 (p. 7): Let 2<a<b be points which are not ordinates, let fin C\textasciicircum{}1[a,b] and G=N-N\_0. Then sum\_a<gamma\_k<= bs\_kf(gamma\_k)=int\_a\textasciicircum{}bG(t)f(t)L(t) dt+(1)/(2)[f G\textasciicircum{}2]\_a\textasciicircum{}b-(1)/(2)int\_a\textasciicircum{}bG(t)\textasciicircum{}2f'(t) dt .}
\expandafter\gdef\csname CGTip@eq:first_moment_identity\endcsname{Eq. (4.7), p. 7: sum\_a<gamma\_k<= bs\_kf(gamma\_k)=int\_a\textasciicircum{}bG(t)f(t)L(t) dt+(1)/(2)[f G\textasciicircum{}2]\_a\textasciicircum{}b-(1)/(2)int\_a\textasciicircum{}bG(t)\textasciicircum{}2f'(t) dt .}
\expandafter\gdef\csname CGTip@cor:cumulative\endcsname{Corollary 4.5 (p. 7): sum\_gamma\_k<= ts\_k=O(log\textasciicircum{}2t). In particular the mean of s\_k over k<= K is O(K\textasciicircum{}-1log\textasciicircum{}2K).}
\expandafter\gdef\csname CGTip@sec:clock_moments\endcsname{Section 5: Moments and extremes of the clock displacement (p. 7)}
\expandafter\gdef\csname CGTip@thm:clock_moments\endcsname{Theorem 5.1 (p. 7): For every fixed integer j>=1, as T->inf, (1)/(no.I\_T)sum\_kin I\_Ts\_k\textasciicircum{}2j=c\_j(log log T)\textasciicircum{}j+O\_j((log log T)\textasciicircum{}j-1/2), c\_j=((2j)!)/(j! (2pi)\textasciicircum{}2j) . In particular the mean square of s\_k over I\_T is (1)/(2pi\textasciicircum{}2)log log T+O(sqrt(log log T)) and (1)/(no.I\_T)sum\_kin I\_Tdelta\_k\textasciicircum{}2=(2log log T)/(log\textasciicircum{}2T)(1+O((log log T)\textasciicircum{}-1/2)).}
\expandafter\gdef\csname CGTip@eq:clock_moments\endcsname{Eq. (5.1), p. 7: (1)/(no.I\_T)sum\_kin I\_Ts\_k\textasciicircum{}2j=c\_j(log log T)\textasciicircum{}j+O\_j((log log T)\textasciicircum{}j-1/2), c\_j=((2j)!)/(j! (2pi)\textasciicircum{}2j) .}
\expandafter\gdef\csname CGTip@eq:delta_mean_square\endcsname{Eq. (5.2), p. 7: (1)/(no.I\_T)sum\_kin I\_Tdelta\_k\textasciicircum{}2=(2log log T)/(log\textasciicircum{}2T)(1+O((log log T)\textasciicircum{}-1/2)).}
\expandafter\gdef\csname CGTip@eq:moments_step1\endcsname{Eq. (5.3), p. 8: sum\_k=m+1\textasciicircum{}n-1P\_j(s\_k)=int\_gamma\_m\textasciicircum{}gamma\_n(S+(omega)/(pi))\textasciicircum{}2jL dt+O(log\textasciicircum{}2j+1T),}
\expandafter\gdef\csname CGTip@eq:moments_step2\endcsname{Eq. (5.4), p. 8: int\_T\textasciicircum{}2TS\textasciicircum{}2jL dt=L(T) c\_jT(log log T)\textasciicircum{}j+O(L(T) T(log log T)\textasciicircum{}j-1/2).}
\expandafter\gdef\csname CGTip@eq:moments_step3\endcsname{Eq. (5.5), p. 8: (1)/(no.I\_T)sum\_kin I\_TP\_j(s\_k)=c\_j(log log T)\textasciicircum{}j+O((log log T)\textasciicircum{}j-1/2).}
\expandafter\gdef\csname CGTip@cor:clock_distribution\endcsname{Corollary 5.2 (p. 8): Let sigma\_T\textasciicircum{}2=(1)/(2pi\textasciicircum{}2)log log T and let Z be a standard normal variable. For every fixed eta >=0, (1)/(no.I\_T)no.kin I\_T:|s\_k|<=eta sigma\_T-> P(|Z|<=eta). In particular, for every fixed A>0, no.kin I\_T:|s\_k|<= A=o(no.I\_T).}
\expandafter\gdef\csname CGTip@eq:clock_distribution\endcsname{Eq. (5.6), p. 8: (1)/(no.I\_T)no.kin I\_T:|s\_k|<=eta sigma\_T-> P(|Z|<=eta).}
\expandafter\gdef\csname CGTip@rem:korolev\endcsname{Remark 5.3 (p. 8): Selberg stated, and Korolev proved [ref], moment formulae for the integer displacement of gamma\_n with respect to the Gram intervals, and Korolev deduced that gamma\_n lies in its own Gram interval only for a set of indices of density zero. By Remark 2.3, Theorem 5.1 and Corollary 5.2 are the analogous statements for...}
\expandafter\gdef\csname CGTip@prop:extremes\endcsname{Proposition 5.4 (p. 8): item limsup\_k->infs\_k=+inf and liminf\_k->infs\_k=-inf. item limsup\_k->inf|delta\_k|<=2pi *0.10076=0.633093.... item On the Riemann hypothesis, delta\_k->0.}
\expandafter\gdef\csname CGTip@eq:sawtooth_sandwich\endcsname{Eq. (5.7), p. 8: s\_k+1-(1)/(2)<= N(t)-N\_0(t)<= s\_k+(1)/(2) .}
\expandafter\gdef\csname CGTip@eq:holder_contradiction\endcsname{Eq. (5.8), p. 8: int\_T\textasciicircum{}2TS\textasciicircum{}2dt<=(int\_T\textasciicircum{}2T|S| dt)\textasciicircum{}2/3(int\_T\textasciicircum{}2TS\textasciicircum{}4dt)\textasciicircum{}1/3<< T(log log T)\textasciicircum{}2/3,}
\expandafter\gdef\csname CGTip@prop:fixed_sup\endcsname{Proposition 5.5 (p. 9): Let X>=2 be fixed and Bin(0,inf]. Then limsup\_k|delta\_k\textasciicircum{}X,B|=limsup\_k|delta\_k|, while limsup\_k|u\_k(X)|=inf.}
\expandafter\gdef\csname CGTip@sec:main_theorem\endcsname{Section 6: The main theorem: a variable cutoff in the truncated equation (p. 9)}
\expandafter\gdef\csname CGTip@subsec:main_statement\endcsname{Section 6.1: Statement (p. 9)}
\expandafter\gdef\csname CGTip@thm:main\endcsname{Theorem 6.1 (p. 9): There are absolute constants T\_0 and C such that for T>= T\_0 and every X with 2<= X<= T\textasciicircum{}1/2, writing c=log X/log T, |(1)/(no.I\_T)sum\_kin I\_Tu\_k(X)\textasciicircum{}2-(1)/(2pi\textasciicircum{}2)log(1)/(c)|<= C(1+sqrt(log(1/c))).}
\expandafter\gdef\csname CGTip@eq:main_theorem\endcsname{Eq. (6.1), p. 9: |(1)/(no.I\_T)sum\_kin I\_Tu\_k(X)\textasciicircum{}2-(1)/(2pi\textasciicircum{}2)log(1)/(c)|<= C(1+sqrt(log(1/c))).}
\expandafter\gdef\csname CGTip@cor:fixed_cutoff\endcsname{Corollary 6.2 (fixed cutoff) (p. 9): Let X>=2 be fixed. Then (1)/(no.I\_T)sum\_kin I\_Tu\_k(X)\textasciicircum{}2=(1)/(2pi\textasciicircum{}2)log log T+O\_X(sqrt(log log T)), and for every fixed A>0, no.kin I\_T:|u\_k(X)|<= A=o(no.I\_T).}
\expandafter\gdef\csname CGTip@eq:fixed_cutoff\endcsname{Eq. (6.2), p. 9: (1)/(no.I\_T)sum\_kin I\_Tu\_k(X)\textasciicircum{}2=(1)/(2pi\textasciicircum{}2)log log T+O\_X(sqrt(log log T)),}
\expandafter\gdef\csname CGTip@cor:power_cutoff\endcsname{Corollary 6.3 (power cutoff) (p. 9): Let 0<c<=(1)/(2) be fixed and X=T\textasciicircum{}c. Then limsup\_T->inf(1)/(no.I\_T)sum\_kin I\_Tu\_k(T\textasciicircum{}c)\textasciicircum{}2<=(1)/(2pi\textasciicircum{}2)log(1)/(c)+C(1+sqrt(log(1/c))), and for every function eta(T)->inf, no.kin I\_T:|u\_k(T\textasciicircum{}c)|>eta(T)=o(no.I\_T). The same holds for X=K\textasciicircum{}c with K=N(T) and fixed 0<c<(1)/(2), because then X=T\textasciicircum{}c\_T with c\_T=c+O(log log T/log T), so...}
\expandafter\gdef\csname CGTip@eq:power_cutoff\endcsname{Eq. (6.3), p. 9: limsup\_T->inf(1)/(no.I\_T)sum\_kin I\_Tu\_k(T\textasciicircum{}c)\textasciicircum{}2<=(1)/(2pi\textasciicircum{}2)log(1)/(c)+C(1+sqrt(log(1/c))),}
\expandafter\gdef\csname CGTip@cor:dichotomy\endcsname{Corollary 6.4 (dichotomy) (p. 9): Let X=X(T) with 2<= X<= T\textasciicircum{}1/2. The mean square of u\_k(X) over kin I\_T is bounded as T->inf if and only if liminf\_T->inflog X/log T>0. In particular, for any fixed X\_0>=2 and fixed cin(0,(1)/(2)], (1)/(no.I\_T)sum\_kin I\_T(u\_k(X\_0)\textasciicircum{}2-u\_k(T\textasciicircum{}c)\textasciicircum{}2)=(1)/(2pi\textasciicircum{}2)log log T+O(sqrt(log log T))->inf .}
\expandafter\gdef\csname CGTip@eq:dichotomy\endcsname{Eq. (6.4), p. 9: (1)/(no.I\_T)sum\_kin I\_T(u\_k(X\_0)\textasciicircum{}2-u\_k(T\textasciicircum{}c)\textasciicircum{}2)=(1)/(2pi\textasciicircum{}2)log log T+O(sqrt(log log T))->inf .}
\expandafter\gdef\csname CGTip@rem:comparison_LeClair\endcsname{Remark 6.5 (p. 9): Corollary 6.4 is the precise sense in which an ordinal-dependent number of primes improves the truncated Franca--LeClair equation. With N=no.p<= X primes, log N=log X+O(log log X), so the condition is that log N grows proportionally to log T. A fixed N, as in the computations of [ref, Sect. VI], leaves a defect whose...}
\expandafter\gdef\csname CGTip@rem:block_rule\endcsname{Remark 6.6 (p. 10): Theorem 6.1 concerns a cutoff that is constant on each window I\_T. Taking T=2\textasciicircum{}r and X=2\textasciicircum{}rc on I\_2\textasciicircum{}r gives a cutoff depending on the ordinal through the window; the rule X\_k=(tau\_k/2pi)\textasciicircum{}theta used in Section 9 differs from such a block rule by a factor at most 2\textasciicircum{}theta inside each window. The extension of (6.1) to cut...}
\expandafter\gdef\csname CGTip@subsec:main_lemmas\endcsname{Section 6.2: Lemmas on prime polynomials (p. 10)}
\expandafter\gdef\csname CGTip@eq:D_def\endcsname{Eq. (6.5), p. 10: D\_1(t)=-(1)/(pi)sum\_pin P(sin(tlog p))/(sqrt(p)), V=sum\_pin P(1)/(p), D=S\textasciicircum{}rm pr\_x\_0-S\_X=epsilon D\_1-S\_X\textasciicircum{}(2)-S\_X\textasciicircum{}(>=3) .}
\expandafter\gdef\csname CGTip@lem:V_size\endcsname{Lemma 6.7 (p. 10): V=log(1/c)+O(1) uniformly for 0<c<=(1)/(2) with X=T\textasciicircum{}c>=2.}
\expandafter\gdef\csname CGTip@lem:trig_MV\endcsname{Lemma 6.8 (p. 10): For real b\_n, 2<= n<= Y, int\_T\textasciicircum{}2T(sum\_n<= Yb\_nsin(tlog n))\textasciicircum{}2dt=(T)/(2)sum\_n<= Yb\_n\textasciicircum{}2+O(sum\_n<= Ynb\_n\textasciicircum{}2+(sum\_n<= Y|b\_n|)\textasciicircum{}2), and the same holds with cos in place of sin.}
\expandafter\gdef\csname CGTip@eq:trig_MV\endcsname{Eq. (6.6), p. 10: int\_T\textasciicircum{}2T(sum\_n<= Yb\_nsin(tlog n))\textasciicircum{}2dt=(T)/(2)sum\_n<= Yb\_n\textasciicircum{}2+O(sum\_n<= Ynb\_n\textasciicircum{}2+(sum\_n<= Y|b\_n|)\textasciicircum{}2),}
\expandafter\gdef\csname CGTip@lem:D_moments\endcsname{Lemma 6.9 (p. 10): With absolute implied constants, int\_T\textasciicircum{}2TD\textasciicircum{}2 dt =(T)/(2pi\textasciicircum{}2)V+O(T(1+sqrt(V))), ; int\_T\textasciicircum{}2TD\textasciicircum{}4 dt << T(1+V)\textasciicircum{}2, ; int\_T\textasciicircum{}2TS\_X'(t)\textasciicircum{}2 dt << T(1+log X)\textasciicircum{}2 .}
\expandafter\gdef\csname CGTip@eq:D_second\endcsname{Eq. (6.7), p. 10: int\_T\textasciicircum{}2TD\textasciicircum{}2 dt =(T)/(2pi\textasciicircum{}2)V+O(T(1+sqrt(V))),}
\expandafter\gdef\csname CGTip@eq:D_fourth\endcsname{Eq. (6.8), p. 10: int\_T\textasciicircum{}2TD\textasciicircum{}4 dt << T(1+V)\textasciicircum{}2,}
\expandafter\gdef\csname CGTip@eq:SX_derivative\endcsname{Eq. (6.9), p. 10: int\_T\textasciicircum{}2TS\_X'(t)\textasciicircum{}2 dt << T(1+log X)\textasciicircum{}2 .}
\expandafter\gdef\csname CGTip@subsec:main_proof\endcsname{Section 6.3: Proof of the main theorem (p. 11)}
\expandafter\gdef\csname CGTip@eq:E0_bounds\endcsname{Eq. (6.10), p. 11: int\_T\textasciicircum{}2TE\_0\textasciicircum{}4 dt<< T,int\_T\textasciicircum{}2TE\_0\textasciicircum{}2 dt<= T\textasciicircum{}1/2(int\_T\textasciicircum{}2TE\_0\textasciicircum{}4 dt)\textasciicircum{}1/2<< T .}
\expandafter\gdef\csname CGTip@eq:J_asymptotic\endcsname{Eq. (6.11), p. 11: J=(T)/(2pi\textasciicircum{}2)V+O(T(1+sqrt(V))),}
\expandafter\gdef\csname CGTip@eq:R_fourth\endcsname{Eq. (6.12), p. 11: int\_T\textasciicircum{}2TR\textasciicircum{}4 dt<< T(1+V)\textasciicircum{}2 .}
\expandafter\gdef\csname CGTip@eq:main_step1\endcsname{Eq. (6.13), p. 11: sum\_kin I\_Tu\_k(X)\textasciicircum{}2=int\_T\textasciicircum{}2TR\textasciicircum{}2L dt+int\_T\textasciicircum{}2TR\textasciicircum{}2S\_X' dt-(no.I\_T)/(12)+O(T\textasciicircum{}3/4).}
\expandafter\gdef\csname CGTip@eq:main_step2\endcsname{Eq. (6.14), p. 11: |int\_T\textasciicircum{}2TR\textasciicircum{}2S\_X' dt|<=(int\_T\textasciicircum{}2TR\textasciicircum{}4)\textasciicircum{}1/2(int\_T\textasciicircum{}2TS\_X'\textasciicircum{}2)\textasciicircum{}1/2<< T(1+V)(1+log X).}
\expandafter\gdef\csname CGTip@eq:main_step3\endcsname{Eq. (6.15), p. 11: (1)/(no.I\_T)sum\_kin I\_Tu\_k(X)\textasciicircum{}2=(V)/(2pi\textasciicircum{}2)+O(1+sqrt(V)+(1+V)/(log T)+(1+V) c),}
\expandafter\gdef\csname CGTip@rem:alternative_route\endcsname{Remark 6.10 (p. 11): An alternative route to an upper bound of the form sum\_kin I\_Tu\_k(T\textasciicircum{}c)\textasciicircum{}2<<\_cno.I\_T averages a pointwise inequality for |u\_k| over windows [gamma\_k,gamma\_k+h] with h=1/log T, and controls the overlap of the windows through the multiplicity M\_h(t)=no.k:gamma\_k<= t<=gamma\_k+h<=1+N(t)-N(t-h), whose moments are bounded by th...}
\expandafter\gdef\csname CGTip@subsec:localization\endcsname{Section 6.4: Localization of the zeros for small cutoffs (p. 12)}
\expandafter\gdef\csname CGTip@prop:localization\endcsname{Proposition 6.11 (p. 12): Let X>=2 with M\_X<=(1)/(2)L(T). Then F\_X is strictly increasing on [T,2T], for every kin I\_T with F\_X(T)<= k-(1)/(2)<= F\_X(2T) the truncated equation F\_X(t)=k-(1)/(2) has exactly one solution t\_k\textasciicircum{}Xin[T,2T], and |gamma\_k-t\_k\textasciicircum{}X|<=(2 |u\_k(X)|)/(L(T)) .}
\expandafter\gdef\csname CGTip@eq:localization\endcsname{Eq. (6.16), p. 12: |gamma\_k-t\_k\textasciicircum{}X|<=(2 |u\_k(X)|)/(L(T)) .}
\expandafter\gdef\csname CGTip@rem:localization_scope\endcsname{Remark 6.12 (p. 12): Since M\_Xasympsqrt(X), Proposition 6.11 applies only for X<<log\textasciicircum{}2T, where log(1/c)\textasciitilde{}log log T; it then turns (6.1) into a mean-square localization of gamma\_k around the solutions of LeClair's truncated equation, but with a growing constant. For X=T\textasciicircum{}c the bound M\_X<=(1)/(2)L(T) fails, F\_X need not be monotone and the...}
\expandafter\gdef\csname CGTip@sec:conditional\endcsname{Section 7: Results conditional on a Landau-type hypothesis (p. 12)}
\expandafter\gdef\csname CGTip@hyp:L\endcsname{Hypothesis 7.1 (L) (p. 12): For every fixed integer n>=2, as t->inf, sum\_0<gamma <= tn\textasciicircum{}igamma=-(t)/(2pi)(Lambda(n))/(sqrt(n))+o(t), the zeros being counted with multiplicity.}
\expandafter\gdef\csname CGTip@eq:hyp_L\endcsname{Eq. (7.1), p. 12: sum\_0<gamma <= tn\textasciicircum{}igamma=-(t)/(2pi)(Lambda(n))/(sqrt(n))+o(t),}
\expandafter\gdef\csname CGTip@prop:RH_implies_L\endcsname{Proposition 7.2 (p. 12): The Riemann hypothesis implies Hypothesis 7.1, with error O(log t).}
\expandafter\gdef\csname CGTip@rem:L_unconditional\endcsname{Remark 7.3 (p. 12): Unconditionally, Theorem 3.8 holds for sum n\textasciicircum{}rho, and |n\textasciicircum{}rho-sqrt(n) n\textasciicircum{}igamma|<= nlog n |beta-(1)/(2)|. Hypothesis 7.1 therefore follows from sum\_0<gamma <= t|beta-(1)/(2)|=o(t), whereas the density estimates of Selberg give O(t) [ref, S 9.24]. Theorems 4.1--6.1 do not use Hypothesis 7.1.}
\expandafter\gdef\csname CGTip@lem:clock_equidistribution\endcsname{Lemma 7.4 (p. 12): For fixed beta !=0, sum\_kin I\_Te\textasciicircum{}ibeta tau\_k<<\_beta log T. Consequently, for fixed integers m,n>=2, (1)/(no.I\_T)sum\_kin I\_Tsin(tau\_klog m)sin(tau\_klog n)->(1)/(2) 1\_m=n.}
\expandafter\gdef\csname CGTip@eq:clock_equidistribution\endcsname{Eq. (7.2), p. 12: (1)/(no.I\_T)sum\_kin I\_Tsin(tau\_klog m)sin(tau\_klog n)->(1)/(2) 1\_m=n.}
\expandafter\gdef\csname CGTip@thm:projection\endcsname{Theorem 7.5 (p. 12): Assume Hypothesis 7.1. For every fixed integer n>=2, (1)/(no.I\_T)sum\_kin I\_Ts\_ksin(gamma\_klog n)->(a\_n)/(2)=-(Lambda(n))/(2pisqrt(n)log n), and the same holds with tau\_k in place of gamma\_k.}
\expandafter\gdef\csname CGTip@eq:projection\endcsname{Eq. (7.3), p. 12: (1)/(no.I\_T)sum\_kin I\_Ts\_ksin(gamma\_klog n)->(a\_n)/(2)=-(Lambda(n))/(2pisqrt(n)log n),}
\expandafter\gdef\csname CGTip@eq:projection_Phi\endcsname{Eq. (7.4), p. 12: Phi(t)=[Sg]\_14\textasciicircum{}t-(1)/(ialpha)sum\_14<gamma <= tn\textasciicircum{}igamma+int\_14\textasciicircum{}tg(L+(omega')/(pi))du=(i)/(alpha)sum\_0<gamma <= tn\textasciicircum{}igamma+O(log t),}
\expandafter\gdef\csname CGTip@eq:projection_partial\endcsname{Eq. (7.5), p. 13: int\_T\textasciicircum{}2TS(t)n\textasciicircum{}itL(t) dt=[Phi L ]\_T\textasciicircum{}2T-int\_T\textasciicircum{}2T(Phi(t))/(2pi t) dt=-i (Lambda(n))/(2pisqrt(n) alpha)int\_T\textasciicircum{}2TL dt+o(Tlog T).}
\expandafter\gdef\csname CGTip@cor:gain\endcsname{Corollary 7.6 (p. 13): Assume Hypothesis 7.1. Let A be a finite set of prime powers and lambda in R. Then (1)/(no.I\_T)sum\_kin I\_T(s\_k\textasciicircum{}2-(s\_k-lambda S\_A(tau\_k))\textasciicircum{}2)->(2lambda-lambda\textasciicircum{}2) v\_A , with v\_A as in (2.11). The limit is maximal for lambda=1, where it equals v\_A.}
\expandafter\gdef\csname CGTip@eq:gain\endcsname{Eq. (7.6), p. 13: (1)/(no.I\_T)sum\_kin I\_T(s\_k\textasciicircum{}2-(s\_k-lambda S\_A(tau\_k))\textasciicircum{}2)->(2lambda-lambda\textasciicircum{}2) v\_A ,}
\expandafter\gdef\csname CGTip@rem:gain_full_factors\endcsname{Remark 7.7 (p. 13): For a finite set of primes P and A\_P,M=p\textasciicircum{}m:pin P,m<= M one has v\_A\_P,Mincreasing to c\_P as M->inf by (2.13). The experiments of Section 9 use the full local factors (M=inf); the passage M->inf inside the limit (7.6) is not covered by the proof, because the convergence in Theorem 7.5 is not shown to be uniform in n....}
\expandafter\gdef\csname CGTip@rem:gain_meaning\endcsname{Remark 7.8 (p. 13): Corollary 7.6 shows that, among the multiples lambda S\_A of the Euler phase, the Euler phase itself (lambda=1) is optimal for the clock, and that its gain v\_A is half of what a perfect predictor of the corresponding Fourier components would gain. Since v\_A\_X->inf as X->inf while v\_A is finite for every fixed A, a fi...}
\expandafter\gdef\csname CGTip@sec:one_step_clock\endcsname{Section 8: The one-step corrected clock (p. 13)}
\expandafter\gdef\csname CGTip@prop:one_step_decomposition\endcsname{Proposition 8.1 (p. 13): Let r\_k(X)=L(tau\_k) delta\_k\textasciicircum{}X be the residual of the one-step clock in ordinal units. Then r\_k(X)=u\_k(X)+(S\_X(gamma\_k)-S\_X(tau\_k))+eps\_k,|eps\_k|<=(delta\_k\textasciicircum{}2)/(2pi min(gamma\_k,tau\_k)), and |S\_X(gamma\_k)-S\_X(tau\_k)|<= M\_X|delta\_k|.}
\expandafter\gdef\csname CGTip@eq:one_step_decomposition\endcsname{Eq. (8.1), p. 13: r\_k(X)=u\_k(X)+(S\_X(gamma\_k)-S\_X(tau\_k))+eps\_k,|eps\_k|<=(delta\_k\textasciicircum{}2)/(2pi min(gamma\_k,tau\_k)),}
\expandafter\gdef\csname CGTip@heur:exponent\endcsname{Heuristic 8.2 (p. 13): This is a heuristic argument. The middle term of (8.1) is approximately S\_X'(gamma\_k) s\_k/L(gamma\_k). Treating S\_X', s\_k and S-S\_X at the zeros as independent centred variables with their t-mean variances, one formally arrives at < r\_k(X)\textasciicircum{}2>\textasciitilde{}(1)/(2pi\textasciicircum{}2)log(1)/(theta)+kappa+theta\textasciicircum{}2 < s\_k\textasciicircum{}2>,theta=(log X)/(log(T/2pi))...}
\expandafter\gdef\csname CGTip@eq:heuristic_rk\endcsname{Eq. (8.2), p. 13: < r\_k(X)\textasciicircum{}2>\textasciitilde{}(1)/(2pi\textasciicircum{}2)log(1)/(theta)+kappa+theta\textasciicircum{}2 < s\_k\textasciicircum{}2>,theta=(log X)/(log(T/2pi)),}
\expandafter\gdef\csname CGTip@eq:heuristic_theta\endcsname{Eq. (8.3), p. 14: theta\_rm opt\textasciitilde{}(1)/(2pisqrt(< s\_k\textasciicircum{}2>)),}
\expandafter\gdef\csname CGTip@rem:negative_exponent\endcsname{Remark 8.3 (p. 14): The following quantitative use of Heuristic 8.2 appears not to work. Inserting the measured < s\_k\textasciicircum{}2> of Table 3 into (8.3) predicts theta\_rm opt=0.614, 0.545, 0.508, 0.488, 0.479, 0.473 in the six windows, that is X\_rm opt\textasciitilde{}20, 42, 90, 158, 214, 270. The observed optima (Table 2) are 7, 43, 151, 401, 701, 907, with l...}
\expandafter\gdef\csname CGTip@heur:second_order\endcsname{Heuristic 8.4 (p. 14): This is a heuristic argument. In Theorem 7.5 the term -(1)/(2)int G\textasciicircum{}2f' of (4.7) and the replacement of gamma\_k by tau\_k are of relative size 1/log T. Modelling S by sum a\_msin(tlog m), the coefficient of cos(tlog n) in S\textasciicircum{}2 for n=p\textasciicircum{}aq\textasciicircum{}b with distinct primes p,q is -a\_p\textasciicircum{}aa\_q\textasciicircum{}b, and the two corrections combine to the...}
\expandafter\gdef\csname CGTip@eq:second_order\endcsname{Eq. (8.4), p. 14: rho\_n:=-2pisqrt(n) < s\_ksin(tau\_klog n)>\textasciitilde{}(Lambda(n))/(log n)+(log n)/(ab log(T/2pi)) 1\_n=p\textasciicircum{}aq\textasciicircum{}b,}
\expandafter\gdef\csname CGTip@sec:numerical_experiments\endcsname{Section 9: Numerical experiments (p. 14)}
\expandafter\gdef\csname CGTip@subsec:data\endcsname{Section 9.1: Data and conventions (p. 14)}
\expandafter\gdef\csname CGTip@exp:data_integrity\endcsname{Experiment 9.1 (p. 14): Object: agreement of the data with independent evaluations. Range: kin 1,10\textasciicircum{}3,10\textasciicircum{}5,10\textasciicircum{}6,K,258317,258318,258319 for the zeros and seven indices for the clock. Expected outcome: agreement to the precision of the data. Observed outcome: the differences between the stored ordinates and mpmath.zetazero are at most 3.5*10...}
\expandafter\gdef\csname CGTip@subsec:methods\endcsname{Section 9.2: Comparison of methods (p. 14)}
\expandafter\gdef\csname CGTip@exp:methods\endcsname{Experiment 9.2 (p. 14): Object: root mean square (RMS), maximum of |residual|, and number of residuals with |residual|>=(1)/(2), for the bare clock, two fixed prime sets with saturation B=0.6, a fixed set of 303 primes without saturation, and the ordinal-dependent cutoff X\_k=max(7,(tau\_k/2pi)\textasciicircum{}0.57) with B=0.6. Range: 1<= k<= K, by windows....}
\expandafter\gdef\csname CGTip@tab:methods\endcsname{Table 1 (p. 15): Residual of the bare clock and of the one-step corrected clocks, by windows of ordinals (Experiment 9.2).}
\expandafter\gdef\csname CGTip@fig:displacement\endcsname{Figure 1 (p. 16): The displacement delta\_k=tau\_k-gamma\_k and the residual of the one-step clock with P=2,3,5,7 and B=0.6, for k<=2000.}
\expandafter\gdef\csname CGTip@fig:rms_height\endcsname{Figure 2 (p. 16): RMS of the residual in logarithmic bins of k for the five methods of Table 1.}
\expandafter\gdef\csname CGTip@fig:max_window\endcsname{Figure 3 (p. 17): Maximum of |residual| in each window for the five methods of Table 1; the dotted line is (1)/(2).}
\expandafter\gdef\csname CGTip@subsec:optimal_cutoff\endcsname{Section 9.3: The optimal cutoff grows with the height (p. 15)}
\expandafter\gdef\csname CGTip@exp:optimal_cutoff\endcsname{Experiment 9.3 (p. 15): Object: RMS of delta\_k\textasciicircum{}X (no saturation) as a function of the cutoff X, for 40 cutoffs between 2 and 1999, and the minimizing cutoff in each window. Range: the six windows. Expected outcome: by Corollary 6.4 the cutoff should grow with the height; Heuristic 8.2 predicts an exponent theta slowly decreasing with the h...}
\expandafter\gdef\csname CGTip@tab:optimal_cutoff\endcsname{Table 2 (p. 17): RMS-optimal prime cutoff of the one-step clock (no saturation) in each window (Experiment 9.3).}
\expandafter\gdef\csname CGTip@fig:rms_cutoff\endcsname{Figure 4 (p. 18): RMS of the one-step residual against the cutoff X (primes p<= X, full local factors), by window; dots mark the minimum of each window.}
\expandafter\gdef\csname CGTip@fig:optimal_cutoff\endcsname{Figure 5 (p. 18): Logarithm of the optimal cutoff against the median of log(tau\_k/2pi) in each window, with the least-squares line through the origin fitted on windows 2--6.}
\expandafter\gdef\csname CGTip@exp:saturation\endcsname{Experiment 9.4 (p. 16): Object: global maximum of |delta\_k\textasciicircum{}X,B|, its location, the number of residuals >=(1)/(2), the RMS, and the residual at k=258318, for the fixed cutoffs 11<= X<=71 and Bin 0.5,0.6,0.7. Range: 1<= k<= K. Expected outcome: none from the theory, which says nothing about finite maxima (Proposition 5.5 concerns the limit)....}
\expandafter\gdef\csname CGTip@subsec:gain_projection\endcsname{Section 9.4: Mean-square gain and projections (p. 17)}
\expandafter\gdef\csname CGTip@exp:gain\endcsname{Experiment 9.5 (p. 17): Object: the gain < s\_k\textasciicircum{}2>-<(s\_k-S\_P(tau\_k))\textasciicircum{}2> for P=2,3,5,7 and P=p<=61 with full local factors, compared with c\_P of (2.13). Range: the six windows. Expected outcome: by Corollary 7.6 and Remark 7.7 (conditional on Hypothesis 7.1), the gain of each finite truncation tends to its v\_A, and v\_Aincreasing to c\_P. Obse...}
\expandafter\gdef\csname CGTip@tab:gain\endcsname{Table 3 (p. 19): Mean-square gain of a fixed Euler phase against its limit c\_P; P\_7=2,3,5,7, P\_61=p<=61 (Experiment 9.5).}
\expandafter\gdef\csname CGTip@fig:gain\endcsname{Figure 6 (p. 19): Measured mean-square gain of a fixed Euler phase in each window (markers) and the constants c\_P (dashed).}
\expandafter\gdef\csname CGTip@exp:projection\endcsname{Experiment 9.6 (p. 17): Object: the normalized projections rho\_n=-2pisqrt(n) < s\_ksin(tau\_klog n)> for 2<= n<=64. Range: 10\textasciicircum{}6<k<= K, where the median of log(tau\_k/2pi) is 11.8382. Expected outcome: rho\_n->Lambda(n)/log n by Theorem 7.5 (conditional on Hypothesis 7.1), with the second-order correction (8.4) at finite height (Heuristic 8.4)....}
\expandafter\gdef\csname CGTip@fig:projection\endcsname{Figure 7 (p. 20): Normalized projections rho\_n of the displacement on sin(tau\_klog n), 10\textasciicircum{}6<k<= K, with the first-order limit Lambda(n)/log n and the second-order prediction (8.4).}
\expandafter\gdef\csname CGTip@subsec:close_pair\endcsname{Section 9.5: The extreme residual: a close pair of zeros (p. 19)}
\expandafter\gdef\csname CGTip@exp:close_pair\endcsname{Experiment 9.7 (p. 19): Object: the zero at which the fixed set P=2,3,5,7 with B=0.6 attains its largest residual. Range: 258315<= k<=258321. Expected outcome: a large |s\_k|, which by (5.7) forces a large value of |S| nearby. Observed outcome: Table 4 and Figure 8. The ordinates gamma\_258317=175681.649131916 and gamma\_258318=175681.7158944...}
\expandafter\gdef\csname CGTip@tab:close_pair\endcsname{Table 4 (p. 20): Ordinates, clock and gaps around the close pair k=258317,258318; the normalized gap is (gamma\_k+1-gamma\_k)log(gamma\_k/2pi)/2pi (Experiment 9.7).}
\expandafter\gdef\csname CGTip@fig:close_pair\endcsname{Figure 8 (p. 20): Zeros, clock and one-step clocks around the close pair k=258317,258318.}
\expandafter\gdef\csname CGTip@sec:discussion\endcsname{Section 10: Discussion (p. 19)}
\expandafter\gdef\csname CGTip@subsec:status\endcsname{Section 10.1: Status of the results (p. 19)}
\expandafter\gdef\csname CGTip@subsec:what_proved\endcsname{Section 10.2: What is proved about the improvement (p. 21)}
\expandafter\gdef\csname CGTip@subsec:novelty\endcsname{Section 10.3: Novelty (p. 21)}
\expandafter\gdef\csname CGTip@rem:no_spacing\endcsname{Remark 10.1 (p. 21): The results proved here concern the arithmetic part of the displacement, of size sqrt(log log T), and its removal by primes; they do not address the local spacing statistics of the zeros. For instance, replacing each gamma\_k by gamma\_k+epsilon\_k with |epsilon\_k|<=epsilon/L(gamma\_k) changes each s\_k by at most epsilo...}
\expandafter\gdef\csname CGTip@sec:open_questions\endcsname{Section 11: Open questions and future work (p. 21)}
\expandafter\gdef\csname CGTip@q:constant\endcsname{Question 11.1 (p. 21): Does (1)/(no.I\_T)sum\_kin I\_Tu\_k(T\textasciicircum{}c)\textasciicircum{}2=(1)/(2pi\textasciicircum{}2)log(1)/(c)+kappa+o(1) hold for fixed cin(0,(1)/(2)], with an explicit constant kappa?}
\expandafter\gdef\csname CGTip@q:large_c\endcsname{Question 11.2 (p. 21): Does Theorem 6.1 extend to (1)/(2)<c<=2, which contains LeClair's range N<[t\textasciicircum{}2]? The proof uses X\textasciicircum{}2<= T in the fourth-moment step of Lemma 6.9.}
\expandafter\gdef\csname CGTip@q:per_ordinate\endcsname{Question 11.3 (p. 21): Does (6.1) hold for cutoffs X\_k=(tau\_k/2pi)\textasciicircum{}theta that vary from zero to zero? This needs a maximal inequality for partial sums of S\_X at the zeros.}
\expandafter\gdef\csname CGTip@q:one_step\endcsname{Question 11.4 (p. 21): Is the mean square of the one-step residual r\_k(T\textasciicircum{}c) asymptotic to c\textasciicircum{}2(2pi\textasciicircum{}2)\textasciicircum{}-1log log T for fixed c, as Heuristic 8.2 suggests, and what is the optimal exponent theta(T) of the one-step clock?}
\expandafter\gdef\csname CGTip@q:experiment_uk\endcsname{Question 11.5 (p. 21): Future experiment: compute u\_k(X) for X=2\textasciicircum{}rc on the windows I\_2\textasciicircum{}r and compare with (6.1); compute the solutions of F\_X(t)=k-(1)/(2) by Newton's method for X<<log\textasciicircum{}2T and compare with Proposition 6.11.}
\expandafter\gdef\csname CGTip@q:generalization\endcsname{Question 11.6 (p. 21): Possible generalization: Dirichlet L-functions, for which Franca and LeClair give the analogous equation [ref] and Selberg's approximation is available in the form of [ref, S 3] for primitive characters.}
\expandafter\gdef\csname CGTip@sec:appendix_tables\endcsname{Section A: Complete tables (p. 21)}
\expandafter\gdef\csname CGTip@tab:scan_rms\endcsname{Table 5 (p. 22): RMS of delta\_k\textasciicircum{}X (no saturation) for all cutoffs of the scan of Experiment 9.3.}
\expandafter\gdef\csname CGTip@tab:scan_max\endcsname{Table 6 (p. 22): Maximum of |delta\_k\textasciicircum{}X| (no saturation) for all cutoffs of the scan of Experiment 9.3.}
\expandafter\gdef\csname CGTip@tab:saturation\endcsname{Table 7 (p. 23): Fixed cutoffs and saturation levels over all K zeros (Experiment 9.4).}
\expandafter\gdef\csname CGTip@tab:projection\endcsname{Table 8 (p. 24): Projections of the displacement, 10\textasciicircum{}6<k<= K (Experiment 9.6). Columns: n; <L(tau\_k)delta\_ksin(tau\_klog n)>; rho\_n; Lambda(n)/log n; second-order prediction log n/(ablog(tau/2pi)) for n=p\textasciicircum{}aq\textasciicircum{}b with the median of log(tau\_k/2pi).}
\expandafter\gdef\csname CGTip@sec:reproducibility\endcsname{Section B: Reproducibility (p. 26)}
\expandafter\gdef\csname CGTip@sec:python_code\endcsname{Section C: Python code (p. 27)}
\makeatother
